\section{Conditional transfer to the matched shell and full block}
\label{sec:full-block}

We now record the strongest valid arithmetic implication.  Every
premise is explicit.

\begin{theorem}[Literal fixed-shift flatness transfer]
\label{thm:literal-flatness-transfer}
Fix \(h_0\ne0\), and let \(a_X\) be the actual determinant coefficient
of the TPC-32 matched packet with
\[
 C_X=\lfloor J_X\rfloor,\qquad
 Q_XJ_X\asymp X,\qquad
 Q_X=X^{267/400+o(1)},\qquad
 J_X=X^{133/400+o(1)}.
\]
Write
\begin{equation}
 N_{0,X}:=J_XQ_X^2\asymp XQ_X
 \label{eq:matched-raw-scale}
\end{equation}
for the literal raw scale of this matched shell.  Suppose there are
fixed exponents \(\eta_Z,\lambda_D\ge0\) for which the following
assumptions hold:
\begin{enumerate}[label=(T\arabic*),leftmargin=2.7em]
\item the complete packet of \cref{def:complete-packet} is literal,
with the same signs, units, orientation, and global raw normalization
as \(a_X\);
\item the missing-Gram hypotheses of
\cref{thm:raw-reassembly} give
\(\delta_{{\rm reasm},X}\);
\item an exact determinant map \(\cJ_X\), a verified subspace \(E_X\),
and a dictionary residual \(r_X\) satisfy
\cref{thm:dictionary-lower};
\item the resulting lower bound proves
\[
 D_X\ge X^{-\lambda_D+o(1)}Q_X^3J_X^2;
 \]
\item independently, the actual zero mode satisfies
\[
 |Z_X|\le X^{-\eta_Z+o(1)}Q_X^2
 \]
with \(2\eta_Z\ge\lambda_D\).
\end{enumerate}
Then
\begin{equation}
 \mathfrak F_{h_0}(X)\le X^{1/400+o(1)}.
 \label{eq:literal-flatness-output}
\end{equation}
Consequently, by the conditional zero-frequency theorem of TPC-32,
every fixed
\begin{equation}
 \sigma<\frac{133}{400}
 \label{eq:matched-saving}
\end{equation}
is a power-saving exponent relative to its raw scale \(N_{0,X}\) for
the complete matched cutoff shell.
Equivalently, if \(\cS_X\) denotes that shell, then
\begin{equation}
 |\cS_X|\le X^{-\sigma+o(1)}N_{0,X}.
 \label{eq:matched-shell-bound}
\end{equation}
\end{theorem}

\begin{proof}
Assumptions (T4)--(T5) and
\cref{cor:zero-slack} give
\eqref{eq:literal-flatness-output}.  Assumption (T1) identifies this
ratio with the distinguished coefficient used in TPC-32, rather than
with a surrogate.  Its conditional high-row closure theorem then
gives \eqref{eq:matched-shell-bound} for every
\(\sigma\) satisfying \eqref{eq:matched-saving}
\citep{WangTPC32}.
\end{proof}

\begin{remark}[What completion contributes]
Assumptions (T2)--(T3) can prove (T4), but they do not prove (T5).
In particular, a complete carrier census, a positive reassembly
constant, and an injective determinant dictionary still do not
cancel the signed zero mode.  The amplitude estimate remains an
arithmetic premise.
\end{remark}

The matched shell is not automatically the whole Type-II residual.
The global step requires an exact block identity of the kind isolated
in TPC-18 \citep{WangTPC18}.

\begin{theorem}[Conditional full-block transfer]
\label{thm:full-block-transfer}
Fix \(h_0\ne0\).  Assume that the physical hard remainder has an exact
decomposition
\begin{equation}
 \cB_{h_0,\delta}(X)
 =
 \sum_{b\in\mathscr B_X}w_{b,X}\cS_{b,X}
 +R_X,
 \qquad
 |\mathscr B_X|=O(\log X),
 \label{eq:full-block-cover}
\end{equation}
where:
\begin{enumerate}[label=(F\arabic*),leftmargin=2.7em]
\item every \(\cS_{b,X}\) is a literal matched packet with complete
native keys and the same fixed shift \(h_0\), and
\(N_{0,b,X}\) denotes its own literal raw scale;
\item the hypotheses of
\cref{thm:literal-flatness-transfer} hold uniformly in \(b\), with
its own literal coefficient \(a_{b,X}\), determinant dictionary, and
certificate data for every block;
\item \(|w_{b,X}|\le X^{o(1)}\), \(R_X=o(X)\), and for some fixed
\(0<\sigma_*<133/400\),
\begin{equation}
 \sum_{b\in\mathscr B_X}
 |w_{b,X}|N_{0,b,X}X^{-\sigma_*}=o(X);
 \label{eq:block-normalization-close}
\end{equation}
\item all calibration, drift, endpoint, nonprimitive, and prime-power
terms omitted from \eqref{eq:full-block-cover} are already included in
\(R_X\) with their literal signs.
\end{enumerate}
Then
\begin{equation}
 \cB_{h_0,\delta}(X)=o(X).
 \label{eq:hard-remainder-small}
\end{equation}
\end{theorem}

\begin{proof}
Choose a fixed
\(\sigma\) with
\(\sigma_*<\sigma<133/400\).
Uniform application of
\cref{thm:literal-flatness-transfer} gives
\[
 |\cS_{b,X}|
 \le X^{-\sigma+o(1)}N_{0,b,X}.
\]
The fixed gap \(\sigma-\sigma_*\) absorbs the uniform \(o(1)\).
Insert the resulting
\(|\cS_{b,X}|\le X^{-\sigma_*}N_{0,b,X}\) in
\eqref{eq:full-block-cover}, sum absolutely, and use
\eqref{eq:block-normalization-close} and \(R_X=o(X)\).
\end{proof}

\begin{remark}[Two different uses of \(\cB\)]
The synthesis operator \(\cB_X=\cA_X+\cM_X\) acts on completion
coefficients.  The scalar \(\cB_{h_0,\delta}(X)\) in
\eqref{eq:full-block-cover} is the physical hard remainder.  The
shared letter does not identify these differently typed objects.
\end{remark}

\begin{remark}[Status of the cover]
\Cref{thm:full-block-transfer} is a logical implication, not a proof
of \eqref{eq:full-block-cover}.  The all-block cover was isolated
earlier as an additional hypothesis.  Nor does the present paper
verify the uniform zero-mode or energy premises for even one growing
physical block.  Therefore \eqref{eq:hard-remainder-small} is not an
unconditional result.
\end{remark}

\begin{remark}[No prime-pair claim]
For an admissible even \(h_0\) with nonzero main-term integral
\(\mathcal I_0(W)\), in the earlier Vaughan normal form a theorem of
the strength
\(\cB_{h_0,\delta}(X)=o(X)\), together with all inherited identities,
would reach a weighted fixed-shift Hardy--Littlewood asymptotic.
That observation measures the strength of the missing premises.  It
does not prove them, and it must not be restated as a twin-prime or
parity result.  For the standard analytic-number-theory framework in
which such reductions sit, see \citet{IwaniecKowalski2004}.
\end{remark}
